\section*{Chapitre \ref{ch:hf:the-business-of-market-making} --- Le métier de teneur de marché}

\begin{solution}{exo:hf:the-business-of-market-making:1}
$2\,145.3/248.5=8.63$ millions de dollars par jour de négociation.
\end{solution}

\begin{solution}{exo:hf:the-business-of-market-making:2}
$10^4\times485.8\times10^6/(1\,940\times10^9)=2.50$ points de base. Le numérateur inclut le revenu de négociation tiré d'autres produits, si bien que le ratio constitue une
borne supérieure sur ce que la société gagnait par euro de produits négociés en bourse.
\end{solution}

\begin{solution}{exo:hf:the-business-of-market-making:3}
$V^\ast=40\,000/(0.3\times10^{-4})=1.33$ milliards de dollars par jour.
\end{solution}

\begin{solution}{exo:hf:the-business-of-market-making:4}
$\pi=0.3\times10^{-4}\times2\times10^9-40\,000=\$20\,000$~; levier $V/(V-V^\ast)=2/(2-1.333)=3.0$. À 1.8~milliards de dollars, $\pi=54\,000-40\,000=
\$14\,000$~: 30~\% de profit en moins pour 10~\% de volume en moins.
\end{solution}

\begin{solution}{exo:hf:the-business-of-market-making:5}
Le mid au moment d'une exécution est périmé~: les exécutions surviennent de façon disproportionnée lorsque le prix efficient a déjà dépassé la cotation (les
traders informés profitent de l'écart, et les cotations mettent du temps à suivre). Face au mid, le Quoter gagne toujours la moitié de la fourchette à l'exécution~; face au prix efficient
elle a déjà perdu 0.173 cent, et le mid rattrape son retard au cours des secondes suivantes, ce qui correspond à la sélection adverse mesurée par rapport au mid.
\end{solution}

\begin{solution}{exo:hf:the-business-of-market-making:6}
Entre 10 secondes ($+0.024$ cent) et 20 secondes ($-0.060$). Toute position conservée plus longtemps qu'environ dix secondes a, en moyenne, perdu sa demi-fourchette~:
le Quoter doit encaisser la fourchette suivante ou se couvrir avant ce délai.
\end{solution}

\begin{solution}{exo:hf:the-business-of-market-making:7}
$-0.150$ cent par action~: le $-0.050$ avec le rabais, moins le rabais de 0.100 cent. Les frais n'entrent pas dans les décisions du Quoter, de sorte que les ordres, les exécutions et
les prix sont identiques~; seule la ligne des frais change.
\end{solution}

\begin{solution}{exo:hf:the-business-of-market-making:8}
Une heure est un tirage. Le P\&L horaire de ce type de Quoter a un écart type voisin de 89~dollars, si bien que 61~dollars se noient largement dans le bruit d'un Quoter dont
la moyenne vraie est négative (celle du chapitre vaut $-\$19$). Décomposer le P\&L, simuler de nombreuses heures indépendantes, et rapporter la moyenne avec son erreur type
avant de dimensionner quoi que ce soit.
\end{solution}

\begin{solution}{pb:hf:the-business-of-market-making:1}
\begin{enumerate}
\item Une société qui fournit de la liquidité par algorithme sur de nombreux instruments et plateformes de négociation~; elle vend de l'immédiateté, tarifée à la moitié de la fourchette.
\item 8.63~millions de dollars par jour~; market making 1\,666 millions de dollars, 77.7~\% (6.71~millions de dollars par jour).
\item Personnel 528.1, communication et données 249.2, opérations 97.9, amortissements 64.4~: 939.6~millions de dollars, soit 3.78~millions de dollars par jour de négociation.
\item 2.09~millions de dollars de revenu de négociation net ajusté et 0.51~millions de dollars de rémunération par salarié.
\item 2.50 points de base~; 0.32~million d'euros de coûts fixes par salarié~; marge d'EBITDA de 41~\%.
\item Revenu de négociation net des frais dépendant du volume, par unité négociée~; $\pi(V)=(c-v)V-F$.
\item $V^\ast=F/(c-v)$~; $d\ln\pi/d\ln V=V/(V-V^\ast)$.
\item 1.33~milliards de dollars~; 20\,000 dollars par jour et un effet de levier de 3.0 à 2~milliards de dollars.
\item $250\times20\,000=5.0$ millions de dollars.
\item $50\,000/(0.3\times10^{-4})=1.67$ milliards de dollars.
\item $40\,000/(0.2\times10^{-4})=2.0$ milliards de dollars~; profit nul à 2~milliards de dollars.
\item Voir la \cref{prop:hf:the-business-of-market-making:split}~: les trois termes entre crochets se télescopent en $m(T)-p_i$.
\item Ses propres ordres déplacent le haut du carnet d'ordres~; s'aligner sur ce haut poursuit ses propres ordres dans une tempête de messages (336\,012 messages contre 323 dans la
simulation du tutoriel).
\item 38\,270 actions par heure, 11~\% du volume du marché, 1\,256 messages.
\item Fourchette $+0.487$, sélection adverse $-0.463$, inventaire $-0.174$, frais $+0.100$~: $-0.050$ cent par action.
\item Ses exécutions se concentrent lorsque le prix efficient a dépassé sa cotation~; le mid au moment de l'exécution est périmé.
\item Écart type de 89~dollars par heure, erreur type $89/\sqrt{10}=\$28$ pour une moyenne de $-\$19$~: indiscernable de zéro.
\item La fourchette rapporte et la position coûte (1.55~euros contre 0.68~euro par transaction)~; le rang de vitesse détermine la performance. Ce Quoter gagne sa fourchette et
le perd sur les positions conservées trop longtemps, et n'a aucun moyen d'être plus rapide que le flux informé.
\item Seuil de rentabilité à 1.33~milliards de dollars par jour~; levier opérationnel de 3.0 à 2~milliards de dollars par jour.
\item Des coûts fixes et un volume de seuil de rentabilité plus bas d'abord~: le taux de capture est fixé par la concurrence et le volume par le marché, tandis qu'une faible marge sur une
grande base de coûts fixes fait de chaque mois calme une perte.
\end{enumerate}
\end{solution}

\begin{solution}{iq:hf:the-business-of-market-making:1}
Immédiateté~: la capacité de négocier tout de suite. Les investisseurs qui franchissent la fourchette la paient, soit la moitié de la fourchette par action~; le coût du teneur de marché est la
sélection adverse venant de ceux qui en savent plus.
\iqlookfor{immédiateté, demi-fourchette, et que le flux informé est le coût.}
\end{solution}

\begin{solution}{iq:hf:the-business-of-market-making:2}
Capture du spread $s\,q\,(m(t)-p)$, sélection adverse $s\,q\,(m(t+H)-m(t))$, inventaire $s\,q\,(m(T)-m(t+H))$. L'horizon sépare le mouvement de prix
que l'exécution prédisait (information) des mouvements ultérieurs de la position détenue (risque)~; le choix de $H$ déplace le P\&L entre les deux termes, jamais le
total.
\iqlookfor{l'identité télescopique et le fait que $H$ est une convention.}
\end{solution}

\begin{solution}{iq:hf:the-business-of-market-making:3}
Sélection adverse~: les exécutions proviennent de traders qui connaissent le prix suivant. Essayer un meilleur prix juste que le mid, pratiquer un skew ou élargir la fourchette lorsque le carnet ou
l'instrument principal signale un mouvement, et il faut se couvrir ou débouclér dans l'horizon sur lequel le mark-out s'estompe.
\iqlookfor{courbes de mark-out, un prix juste, et agir dans l'horizon de décroissance.}
\end{solution}

\begin{solution}{iq:hf:the-business-of-market-making:4}
Une boucle de rétroaction~: le Quoter lit un carnet qui contient ses propres ordres et ne cesse de recalculer son prix pour s'aligner sur lui-même ou s'en écarter. Coter à partir du carnet sans
ses propres ordres, et ajouter une variation minimale de prix avant de recalculer.
\iqlookfor{auto-référence dans la vue des données de marché, et un test de cadence de messages.}
\end{solution}

\begin{solution}{iq:hf:the-business-of-market-making:5}
Coûts fixes~: le profit vaut $(c-v)V-F$, de sorte que sa variation en pourcentage est $V/(V-V^\ast)$ fois celle du volume~; et la capture elle-même varie avec la volatilité et
la concurrence.
\iqlookfor{levier opérationnel et volume de seuil de rentabilité.}
\end{solution}

\begin{solution}{iq:hf:the-business-of-market-making:6}
$2\times89/\sqrt n=19$ donne $n=(178/19)^2\approx88$ heures.
\iqlookfor{loi d'échelle de l'erreur type et un nombre.}
\end{solution}
